\documentclass[11pt]{amsart}
\usepackage{amsmath,amssymb}
\setlength{\emergencystretch}{3em}
\tolerance=600
\usepackage{tikz-cd}
\usepackage{tikz}
\usepackage[hidelinks]{hyperref}

\newtheorem{thm}{Theorem}[section]
\newtheorem{lem}[thm]{Lemma}
\newtheorem{cor}[thm]{Corollary}
\newtheorem{prop}[thm]{Proposition}
\theoremstyle{definition}
\newtheorem{defi}[thm]{Definition}
\newtheorem{rem}[thm]{Remark}
\newtheorem{ex}[thm]{Example}
\newtheorem{que}[thm]{Question}

\newcommand{\Set}{\mathbf{Set}}
\newcommand{\Fin}{\mathbf{FinSet}}
\newcommand{\Sub}{\mathrm{Sub}}
\newcommand{\Mod}{\mathrm{Mod}}
\newcommand{\Sent}{\mathrm{Sent}}
\newcommand{\Lan}{\mathrm{Lan}}
\newcommand{\Aut}{\mathrm{Aut}}
\newcommand{\PP}{\mathsf P}
\newcommand{\HH}{\mathsf H}
\newcommand{\cX}{\mathcal X}
\newcommand{\cC}{\mathcal C}
\newcommand{\cE}{\mathcal E}
\newcommand{\cM}{\mathcal M}
\newcommand{\cS}{\mathcal S}
\newcommand{\cJ}{\mathcal J}
\newcommand{\bT}{\mathbb T}
\newcommand{\id}{\mathrm{id}}

\begin{document}

\title[Finite sorts and finitely many sentences]{Intermediate logics
and the imaginaries of a theory with finite sorts and finitely many
sentences: profinite leg trees}
\author{Ryota Shibusawa}
\thanks{Corresponding author. E-mail: \texttt{r-shibusawa@daiichi-koudai.ac.jp}}
\address{Daiichi Institute of Technology}
\email{r-shibusawa@daiichi-koudai.ac.jp}
\keywords{intermediate logic, G\"odel--Dummett logic, Heyting category, pretopos completion, imaginaries, Kripke model, EI-category, profinite group, coherent theory}
\subjclass[2020]{03G30 (primary), 03B55, 03C13, 18F10, 20E18 (secondary)}

\begin{abstract}
For an intermediate logic $L$ and a coherent theory, reading the
theory in first-order logic over $L$ and adding the imaginaries need
not commute. For the theories all of whose models are finite on every
sort the question was left open for the logics proving linearity which
have a chain of three elements as a frame. We settle it when the
theory has finitely many sentences up to equivalence: the two
operations commute if and only if every endomorphism of a model is invertible and there is no
obstruction by a profinite leg tree. Such a tree is built from chains
of non-invertible homomorphisms of models, and the condition compares
two closed subgroups of the profinite automorphism group of a model.
For a family of theories given by closed subgroups $A,T_1,\dots,T_m$
of a profinite abelian group the condition is $A\cap\sum_{j\in
J}T_j=\sum_{j\in J}(A\cap T_j)$ for all $J$. An example over the
$p$-adic integers shows that the absence of obstructions is not
monotone along finite approximations of a theory. We also prove that
commutation is inherited by the theories obtained by making some
sentences true and others false, and that it holds for all logics
proving linearity when every endomorphism of a model is an identity.
\end{abstract}

\maketitle

\section{Introduction}

For an intermediate logic $L$ and a small coherent category $\cC$ let
$\HH_L\cC$ be the syntactic category of the internal coherent theory of
$\cC$ in first-order logic over $L$, and $\PP$ the pretopos completion,
which adds the imaginaries. There is a full and faithful comparison
$c^L_\cC\colon\HH_L\PP\cC\to\PP\HH_L\cC$, and it is an equivalence if
and only if $\HH_L\PP\cC$ is a pretopos \cite[Lemma~2.2]{Interm}. We
say that \emph{the translation into $L$ commutes with the imaginaries
of $\cC$} if $c^L_\cC$ is an equivalence.

A \emph{model} of $\cC$ is a coherent functor $M\colon\cC\to\Set$, a
\emph{homomorphism} is a natural transformation, and $\cC$ \emph{has
finite sorts} if $M(A)$ is finite for every model $M$ and object $A$.
For these $\cC$ the comparison was determined in \cite{Finsorts} in
terms of the category $\Mod(\cC)$ of models, for all $L$ except those
proving the linearity axiom $(p\to q)\vee(q\to p)$ with $n_L\ge3$; here
$n_L\in\{1,2,\dots,\infty\}$ is the supremum of the numbers of
elements of the finite chains which are frames of $L$. In that case it
is necessary that $\Mod(\cC)$ is an \emph{EI-category} (every
endomorphism is invertible), and this is not sufficient
\cite{Strings}. This paper treats that case.

A \emph{sentence} is a subobject of the terminal object $1$ of $\cC$.
With finite sorts, the models lie over the prime filters of the
lattice $\Sub(1)$ of sentences: there is a homomorphism $M\to N$ iff
the sentences true in $M$ are true in $N$, and in the EI case the
isomorphism classes of models are the prime filters and the
non-invertible homomorphisms are those which enlarge the set of true
sentences (Proposition~\ref{prop:spec}). A \emph{string} is a sequence
$M_0\to\dots\to M_r$ of non-invertible homomorphisms. If $\Sub(1)$ is
finite, the automorphism group $G_s$ of the first model of a string
$s$ is a profinite group, and it contains the closed subgroup $H_s$ of
the automorphisms which extend to an automorphism of the string. A
\emph{profinite leg tree} $\mathcal T$ at $s$ (Section~\ref{sec:trees})
is given by a finite tree of strings which extend $s$ to the left, for
the restrictions of the models to finitely many sorts, together with
all its liftings to strings of models; it
defines closed subgroups $R_{\mathcal T}\le H_s\cap S_{\mathcal T}$ of
$G_s$, by forming joins of equivalence relations in two ways, on all
homomorphisms and on those which are homomorphisms of strings, and it
is an \emph{obstruction} if $R_{\mathcal T}\ne H_s\cap S_{\mathcal T}$.

\begin{thm}\label{thm:A}
Let $\cC$ be a small coherent category with finite sorts and finitely
many sentences, and $L$ a non-classical intermediate logic proving
linearity. Then $c^L_\cC$ is an equivalence if and
only if $\Mod(\cC)$ is an EI-category and no profinite leg tree with
strings of at most $n_L$ models is an obstruction.
\end{thm}

For a finite category of models this is \cite[Thm.~6.15]{Strings},
where all groups are finite and the joins are the ordinary ones. The
first examples with infinite groups are the following. Let $G$ be a
profinite abelian group and $A,T_1,\dots,T_m$ closed subgroups, and
let $\cC(G;A;T_\bullet)$ be the pretopos of the continuous functors to
finite sets on the category
\[
\begin{tikzcd}[column sep=large]
0_j\ar[r,"G/T_j"]\ar[rr,bend right=20,"G/T_j"']&1\ar[r,"G"]&2\ar[loop right,"A"]
\end{tikzcd}\qquad(1\le j\le m)
\]
with composition the addition of $G$ (Section~\ref{sec:three}).

\begin{thm}\label{thm:B}
Let $L$ prove linearity with $n_L\ge3$. Then $c^L_\cC$ is an
equivalence for $\cC=\cC(G;A;T_\bullet)$ if and only if
\begin{equation}\label{eq:dist}
A\cap\sum_{j\in J}T_j=\sum_{j\in J}(A\cap T_j)\qquad\text{for every }
J\subseteq\{1,\dots,m\}.
\end{equation}
In this case the open subgroups $N$ of $G$ for which the finite
quotient $(G/N;(A+N)/N;(T_\bullet+N)/N)$ satisfies \eqref{eq:dist}
form a base of neighbourhoods of $0$.
\end{thm}

For $G=\mathbb Z_p^3$, $T_1=\mathbb Z_pe_1$, $T_2=\mathbb Z_pe_2$ and
$A=\mathbb Z_p(1,1,p)$ the condition holds, while it fails for every
quotient $G/p^kG$ (Example~\ref{ex:Zp}): a theory for which the
comparison is an equivalence can be the union of a chain of theories
with finite categories of models for none of which it is.

Two further results do not need finitely many sentences.

\begin{thm}[localisation]\label{thm:C}
Let $\cC$ be a pretopos, $L$ an intermediate logic, $F$ a filter and
$I$ an ideal of $\Sub(1)$, and $\cC_{F,I}$ the pretopos of the theory
of $\cC$ with the additional axioms $\top\vdash s$, $s\in F$, and
$s'\vdash\bot$, $s'\in I$. If $c^L_\cC$ is an equivalence, so is
$c^L_{\cC_{F,I}}$.
\end{thm}

So for finite sorts and $L$ non-classical proving linearity, every part of
$\Mod(\cC)$ cut out by a filter and an ideal which is essentially
finite must be an EI-category without the obstructions of
\cite{Strings} (Corollary~\ref{cor:nec}).

\begin{thm}[rigid models]\label{thm:D}
Let $\cC$ be a small coherent category with finite sorts such that
every endomorphism of every model is an identity. Then $c^L_\cC$ is an
equivalence for every $L$ proving linearity.
\end{thm}

For propositional theories, whose models form a poset, this is
\cite[Thm.~6.3]{Interm}; for $n_L\le2$ it is contained in
\cite[Thm.~1.1]{Finsorts}.

\subsection*{The method}
A pretopos $\cC$ with finite sorts, EI models and finitely many
sentences is the directed colimit of pretoposes $\Fin^{\cX_S}$, where
$\cX_S$ is the finite category of the restrictions of the models to a
finite set $S$ of sorts; the transition functors $\cX_{S'}\to\cX_S$
are bijective on objects and full, and $\Mod(\cC)$ is the limit
(Proposition~\ref{prop:general}).
\[
\begin{tikzcd}[column sep=large,row sep=small]
\Mod(\cC)\simeq\varprojlim\cX_S\ar[r]&\cX_{S'}\ar[r,"q"]&\cX_S\\
\cC=\varinjlim\Fin^{\cX_S}&\Fin^{\cX_{S'}}\ar[l]&\Fin^{\cX_S}\ar[l,"q^*"']
\end{tikzcd}
\]
By \cite{Strings}, $\PP\HH_L\Fin^{\cX}$ is the category of the functors
$W$ from the strings of $\cX$ to finite sets, and $\HH_L\Fin^\cX$
consists of those $W$ for which a unit map is injective. The
comparison for $\cC$ is an equivalence iff every $W$ at a stage $S$
has an injective unit at some later stage (Lemma~\ref{lem:eventual}).
The kernel of the unit of $W$ is controlled by the leg trees which are
realised in $W$. The point is that among these there is a largest one
for each element of $W$, and that the largest tree for $W\circ q$ is
the pullback of the largest tree for $W$
(Lemma~\ref{lem:mono} and the proof of Theorem~\ref{thm:proftree}); so the trees at the different stages form
inverse systems and compactness applies.

\subsection*{What is new and what is known}
Known: the classification of \cite{Finsorts}, with the necessity of
the EI property; the criterion for a finite category of models
\cite[Thm.~1.1, Thm.~6.15]{Strings} and its case
\cite[Prop.~6.11]{Strings} for finite abelian groups;
Theorem~\ref{thm:D} for propositional theories \cite{Interm} and for
finite categories of models \cite[Cor.~6.5]{Strings}. New, to the
author's knowledge: Theorems~\ref{thm:A}--\ref{thm:D} beyond these
cases. Proposition~\ref{prop:finmod}, that a pretopos with finite
sorts and an essentially finite category of models is $\Fin^\cX$, is
presumably known; see \cite[Cor.~2.6]{Ye24}. We include a short proof
from the lemma on types of \cite{Finsorts}. A duality between
profinite monoids and the pretoposes of the categorical coherent
theories is given in \cite{Ye26}; we do not use it.

\section{The setting and the results used}\label{sec:setting}

Conventions are those of \cite{Finsorts}. All coherent categories are
small. An \emph{$L$-category} is a Heyting category whose lattices of
subobjects are Heyting algebras in which the formulas of $L$ are
valid; $\eta\colon\cC\to\HH_L\cC$ is the free $L$-category on $\cC$:
for an $L$-category $\cE$, composition with $\eta$ is an equivalence
from the groupoid of the Heyting functors $\HH_L\cC\to\cE$ and natural
isomorphisms to that of the coherent functors $\cC\to\cE$. We omit
$\eta$. The models of $\PP\cC$ are those of $\cC$, $\PP\cC$ has finite
sorts and the same sentences if $\cC$ has, and $c^L_\cC$ is an
equivalence iff $c^L_{\PP\cC}$ is. So \emph{from now on $\cC$ is a
pretopos}; then $c^L_\cC$ is an equivalence iff $\HH_L\cC$ is a
pretopos \cite[Lemma~2.2]{Interm}. For a model $M$ let $\Sent(M)$ be
the prime filter of the sentences true in $M$. Directed colimits of
categories are taken along functors compatible up to coherent
isomorphisms. We use:

\begin{itemize}
\item[(R1)] (Completeness.) $U\le V$ for subobjects of $A$ iff
$U(M)\subseteq V(M)$ for all models $M$; every prime filter of
$\Sub(1)$ is $\Sent(M)$ for some $M$.
\item[(R2)] (Types, \cite[Lemma~4.3]{Finsorts}.) Let $N$, $N'$ be
models of a coherent category which are finite on every object, and
$(\alpha_i)_{i\in I}$, $(\alpha'_i)_{i\in I}$ families of elements
$\alpha_i\in A_i(N)$, $\alpha'_i\in A_i(N')$ such that for all
$i_1,\dots,i_r\in I$, $r\ge0$, and all $S\le A_{i_1}\times\dots\times
A_{i_r}$: $(\alpha_{i_1},\dots,\alpha_{i_r})\in S(N)$ implies
$(\alpha'_{i_1},\dots,\alpha'_{i_r})\in S(N')$. Then there is a
homomorphism $\lambda\colon N\to N'$ with $\lambda(\alpha_i)=
\alpha'_i$ for all $i$.
\item[(R3)] (\cite[Lemma~5.3 and its proof]{Finsorts}.) If a pretopos
$\cC$ is the colimit of a directed system of pretoposes $\cC_i$ and
coherent functors, then $\HH_L\cC$ is the colimit of the $\HH_L\cC_i$;
and if $\HH_L\cC_i$ is a pretopos for all $i$ in a cofinal subset,
$\HH_L\cC$ is a pretopos.
\item[(R4)] (\cite[Remark after Thm.~5.5]{Finsorts}.) If $\cC$ has
finite sorts, $L$ is not classical and $c^L_\cC$ is an equivalence,
then $\Mod(\cC)$ is an EI-category.
\end{itemize}

\subsection*{Finite categories of models}
Let $\cX$ be a finite EI-category, $L$ a logic proving linearity and
$n=n_L$. The models of $\Fin^\cX$ are the evaluations at the objects
and the homomorphisms are the arrows of $\cX$. A \emph{string} is a
sequence $s=(m_0\to\dots\to m_r)$ of non-invertible arrows, $r\ge0$,
with at most $n$ objects (for $n=\infty$: of any length; the length is
at most the number of isomorphism classes of objects); its final segments are the $s_{\ge
j}=(m_j\to\dots\to m_r)$, and an isomorphism of strings is a family of
isomorphisms of the objects commuting with the arrows. The category
$\cS_n(\cX)$ has the strings as objects, an arrow $s\to t$ being an
arrow $g\colon s_0\to t_0$ of $\cX$ of the form $\iota_0p$ with $p$
the composite $s_0\to s_j$ of arrows of $s$ and $\iota$ an isomorphism
of strings from $s_{\ge j}$ onto $t$; the index $j$ is determined by
$t$. Let $\pi\colon\cS_n(\cX)\to\cX$ be the first object. (In
\cite{Strings} one string is taken in each isomorphism class; this
gives an equivalent category.)

\begin{itemize}
\item[(R5)] (\cite[Thm.~1.1]{Strings}.) $\PP\HH_L\Fin^\cX\simeq
\Fin^{\cS_n(\cX)}$, by the functor extending $U\mapsto(s\mapsto\hat
G_s(U))$, where $\hat G_s$ is the model of $\HH_L\Fin^\cX$ at the
first node of the Kripke model $s$ over a chain; $\HH_L\Fin^\cX$
corresponds to the full subcategory of the subfunctors of the functors
$\pi^*A$, $A$ in $\Fin^\cX$; and $\HH_L\Fin^\cX$ is a pretopos iff for
every $W\colon\cS_n(\cX)\to\Fin$ the unit $W\to\pi^*\Lan_\pi W$ is
injective. A functor $W$ is isomorphic to a subfunctor of some
$\pi^*A$ iff its unit is injective.
\end{itemize}

For strings $t$, $s$ let $V_t=\cX(t_0,s_0)$ and $K_t=\cS_n(t,s)
\subseteq V_t$; $t$ is \emph{over} $s$ if $K_t\ne\emptyset$. A
\emph{one-step arrow} is an arrow $\beta\colon t'\to t$ of $\cS_n$
where $t'$ has one object more than $t$; if $t$ is over $s$, then $t'$
is over $s$ and $K_{t'}=K_t\beta$. A \emph{leg tree} $T$ at $t$
\cite[Def.~6.14]{Strings} is a finite rooted tree whose root is
labelled by $t$ and whose other nodes are labelled by one-step arrows
into the string of the parent node; so the root has children
$(\beta,T')$ with $\beta\colon t'\to t$ one-step and $T'$ a leg tree
at $t'$. Equivalence relations $E_T$ on $V_t$ and $F_T$ on $K_t$ are
defined by recursion: the equality if the root has no children, and
otherwise
\[
E_T=\bigvee\beta^*E_{T'},\qquad F_T=\bigvee\beta^*F_{T'},
\]
the joins of equivalence relations over the children of the root,
where $g\mathrel{\beta^*E}g'$ iff $g\beta\mathrel Eg'\beta$. A
\emph{realisation} of $T$ in $W\colon\cS_n\to\Fin$ is a family
$y_\nu\in W(t_\nu)$, $\nu$ a node with string $t_\nu$, with
$W(\beta)y_\nu=y_{\nu'}$ for $\nu$ a child of $\nu'$ with label
$\beta$. A \emph{morphism} of leg trees at the same string is a map of
trees preserving roots, parents and labels. For $t=s$ we have
$V_s=G_s:=\Aut_\cX(s_0)$ and $K_s=H_s:=\Aut_{\cS_n}(s)$.
\[
\begin{tikzcd}[column sep=small,row sep=small]
t''_1\ar[dr]&&t''_2\ar[dl]&&\\
&t'_1\ar[drr,"\beta_1"']&&t'_2\ar[d,"\beta_2"]&t'_3\ar[dl,"\beta_3"]\\
&&&t&
\end{tikzcd}
\]

\begin{itemize}
\item[(R6)] (\cite[proof of Thm.~6.15, (a), (b) and the Claim]{Strings}.)
For $W\colon\cS_n(\cX)\to\Fin$, a string $s$ and $a,b\in W(s)$: the
unit of $W$ identifies $a$ and $b$ if and only if there are $\Theta\in
H_s$ and a leg tree $T$ at $s$ realised in $W$ with $a$ at the root
such that $b=W(\Theta)a$ and $\id\mathrel{E_T}\Theta$. If $T$ is
realised in $W$ with $a$ at the root and $\id\mathrel{F_T}\Theta$,
then $W(\Theta)a=a$.
\item[(R7)] (\cite[proof of Thm.~6.15, (i)$\Rightarrow$(ii)]{Strings}.)
For every leg tree $T$ at $s$ there is $W_T\colon\cS_n(\cX)\to\Fin$
with a realisation of $T$, with $a$ at the root, such that for
$\Theta\in H_s$: $W_T(\Theta)a=a$ iff $\id\mathrel{F_T}\Theta$.
\end{itemize}
The relations at the root $s$ are invariant under left multiplication:
$E_T$ is $g^{-1}g'\in S_T$ and $F_T$ is $h^{-1}h'\in R_T$ for
subgroups $R_T\le H_s\cap S_T$ of $G_s$ \cite[\S6]{Strings}. In
\cite{Strings} these statements are made for a category of strings
with one string in each isomorphism class; the units, the leg trees
and the relations are transported along the equivalence with
$\cS_n(\cX)$, an isomorphism of strings $t\cong t_1$ identifying
$V_t$, $K_t$ and the leg trees at $t$ with those at $t_1$.

\section{The models and the prime filters of sentences}

Let $\cC$ be a pretopos with finite sorts.

\begin{prop}\label{prop:spec}
\begin{enumerate}
\item[(i)] There is a homomorphism $M\to N$ iff $\Sent(M)\subseteq
\Sent(N)$.
\item[(ii)] If $\Mod(\cC)$ is an EI-category, a homomorphism $M\to N$
is invertible iff $\Sent(M)=\Sent(N)$; so the isomorphism classes of
models correspond to the prime filters of $\Sub(1)$.
\item[(iii)] If $\Mod(\cC)$ is an EI-category, there are two
composable non-invertible homomorphisms iff $\Sub(1)$ has prime
filters $P_0\subsetneq P_1\subsetneq P_2$.
\end{enumerate}
\end{prop}
\begin{proof}
(i) Homomorphisms preserve the truth of sentences; the converse is
(R2) for the empty family. (ii) If $\Sent(M)=\Sent(N)$ and $h\colon
M\to N$, there is $g\colon N\to M$ by (i), and $gh$, $hg$ are
invertible, so $h$ is; the converse is clear, and every prime filter
is realised by (R1). (iii) follows from (i) and (ii).
\end{proof}

So the conditions of \cite[Thm.~1.1]{Finsorts} read: for the logics
without the weak excluded middle, $\Mod(\cC)$ is an EI-category and
$\Sub(1)$ is a Boolean algebra; for those with the weak excluded
middle and without linearity, $\Mod(\cC)$ is an EI-category and
$\Sub(1)$ has no chain of three prime filters.

\section{Localisation}\label{sec:loc}

Let $\cC$ be a coherent category and $s\le1$. The slice $\cC/s$ is the
full subcategory of the objects $A$ whose map to $1$ factors through
$s$; it is a coherent category with terminal object $s$, a pretopos if
$\cC$ is, and $s^*\colon\cC\to\cC/s$, $A\mapsto A\times s$, is
coherent.

\begin{lem}\label{lem:slice}
\begin{enumerate}
\item[(a)] For a coherent category $\cE$, composition with $s^*$ is an
equivalence between the coherent functors $\cC/s\to\cE$ and the
coherent functors $F\colon\cC\to\cE$ with $F(s)\cong1$.
\item[(b)] If $\cC$ is an $L$-category, so is $\cC/s$, $s^*$ is
Heyting, and (a) holds for $L$-categories $\cE$ and Heyting functors.
\item[(c)] $\HH_L(\cC/s)\simeq(\HH_L\cC)/\eta s$.
\end{enumerate}
\end{lem}
\begin{proof}
(a) The inclusion $\cC/s\to\cC$ preserves pullbacks, binary products
(as $A\times_sB=A\times B$ for a subterminal $s$), images and joins of
subobjects, and $\Sub_{\cC/s}(A)=\Sub_\cC(A)$. So for $F$ with
$F(s)\cong1$ the restriction $F|$ of $F$ to $\cC/s$ is coherent, and
$F|s^*(A)=F(A)\times F(s)\cong F(A)$. Conversely $s^*A\cong A$ for $A$
in $\cC/s$, so a coherent $G\colon\cC/s\to\cE$ is isomorphic to
$(Gs^*)|$, and $Gs^*(s)=G(s)\cong1$.

(b) For $f\colon A\to B$ in $\cC/s$ the maps $f^*$, $\exists_f$,
$\forall_f$ between $\Sub(A)$ and $\Sub(B)$ are those of $\cC$, so
$\cC/s$ is a Heyting category with the same Heyting algebras of
subobjects, an $L$-category. The base change $s^*\colon\cC\to\cC/s$ of a Heyting category
preserves $\forall$ and $\Rightarrow$ by the Beck--Chevalley
condition, so it is Heyting. A Heyting functor $F$ on $\cC$ with $F(s)
\cong1$ restricts to a Heyting functor on $\cC/s$ by the same remark.

(c) For an $L$-category $\cE$: Heyting functors $(\HH_L\cC)/\eta s\to
\cE$ are Heyting functors $G$ on $\HH_L\cC$ with $G(\eta s)\cong1$ by
(b); these are coherent functors $F$ on $\cC$ with $F(s)\cong1$; these
are coherent functors on $\cC/s$ by (a). The corresponding coherent
functor $\cC/s\to(\HH_L\cC)/\eta s$ is the restriction of $\eta$.
\end{proof}

For the closed part let $\cC$ be a pretopos, $s\le1$ and
$\cC_{\neg s}$ the syntactic category of $\cC$ with the additional
axiom $s\vdash\bot$: its objects are those of $\cC$, $\Sub_{\cC_{\neg
s}}(A)$ is the quotient of $\Sub_\cC(A)$ by the congruence $U\equiv V$
iff $U\vee s_A=V\vee s_A$ ($s_A=A\times s$), and the morphisms are the
functional relations modulo this congruence. The quotient functor
$\kappa\colon\cC\to\cC_{\neg s}$ is coherent, surjective on objects and
on subobjects, and coherent functors $\cC_{\neg s}\to\cE$ correspond to
coherent functors $F$ on $\cC$ with $F(s)\cong0$. The models of
$\cC_{\neg s}$ are the models of $\cC$ in which $s$ is false.

\begin{lem}\label{lem:closed}
\begin{enumerate}
\item[(a)] $\cC_{\neg s}$ is a pretopos.
\item[(b)] $\HH_L(\cC_{\neg s})\simeq(\HH_L\cC)/\neg\eta s$.
\end{enumerate}
\end{lem}
\begin{proof}
(a) A coherent functor preserves the property of an object to be the
disjoint join of two subobjects, and in a coherent category such an
object is a coproduct; so $\kappa(A\sqcup B)$ is a disjoint coproduct
of $\kappa A$ and $\kappa B$. An equivalence relation on $\kappa A$ is
$\kappa R$ with $R\le A\times A$, and $R'=R\vee s_{A\times A}$ is an
equivalence relation on $A$ in $\cC$: it is reflexive and symmetric
because $R$ is so modulo the congruence, and $R'\circ R'\le R'$ because
a composite involving $s_{A\times A}$ lies over $s$ and $R\circ R\le
R'$. Its quotient $A\to A/R'$ in $\cC$ is a cover with kernel pair
$R'$, and $\kappa$ preserves covers and kernel pairs; in a regular
category a cover is the coequaliser of its kernel pair. So $\kappa R=
\kappa R'$ is effective.

(b) A Heyting functor $G$ has $G(\eta s)\cong0$ iff $G(\neg\eta s)\cong
1$. So for an $L$-category $\cE$ the Heyting functors $(\HH_L\cC)/\neg
\eta s\to\cE$ are, by Lemma~\ref{lem:slice}(b), the Heyting functors
$G$ on $\HH_L\cC$ with $G(\eta s)\cong0$, that is the coherent $F$ on
$\cC$ with $F(s)\cong0$, that is the coherent functors on $\cC_{\neg
s}$.
\end{proof}

For $s,s'\le1$ we write $(\cC/s)_{\neg s'}$ for the pretopos obtained
from $\cC/s$ by the axiom $s\wedge s'\vdash\bot$, $s\wedge s'$ being a
subobject of the terminal object $s$ of $\cC/s$. By
Lemma~\ref{lem:slice}(a) and the property of $\kappa$, the coherent
functors $(\cC/s)_{\neg s'}\to\cE$ correspond to the coherent functors
$F$ on $\cC$ with $F(s)\cong1$ and $F(s')\cong0$; so for $s_1\le s$ and
$s'\le s'_1$ there is a coherent functor $(\cC/s)_{\neg s'}\to
(\cC/s_1)_{\neg s'_1}$ under $\cC$. For a filter $F$ and an ideal $I$
of $\Sub(1)$ let $\cC_{F,I}$ be the colimit of this system over $s\in
F$ (downwards) and $s'\in I$ (upwards), which is directed. It is a
pretopos; the coherent functors on it correspond to the coherent
functors on $\cC$ which send the sentences in $F$ to $1$ and those in
$I$ to $0$, so it is the pretopos of the theory of $\cC$ with the
axioms $\top\vdash s$, $s\in F$, and $s'\vdash\bot$, $s'\in I$; and its
models are the models $M$ of $\cC$ with $F\subseteq\Sent(M)$ and
$\Sent(M)\cap I=\emptyset$.
\[
\begin{tikzcd}[column sep=huge]
\cC\ar[r,"\eta"]\ar[d,"s^*"']&\HH_L\cC\ar[d,"(\eta s)^*"]\\
\cC/s\ar[r]&\HH_L(\cC/s)\simeq(\HH_L\cC)/\eta s
\end{tikzcd}
\]

\begin{thm}[localisation]\label{thm:loc}
Let $\cC$ be a pretopos such that $\HH_L\cC$ is a pretopos. Then
$\HH_L\cC_{F,I}$ is a pretopos for every filter $F$ and ideal $I$ of
$\Sub(1)$.
\end{thm}
\begin{proof}
Slices of pretoposes are pretoposes, so $\HH_L((\cC/s)_{\neg s'})$ is a
pretopos by Lemmas~\ref{lem:slice}(c) and~\ref{lem:closed}; then (R3).
\end{proof}


\begin{proof}[Proof of Theorem~\ref{thm:C}]
This is Theorem~\ref{thm:loc}: the theory of $\cC$ with the additional
axioms is that of $\cC_{F,I}$.
\end{proof}

\section{Finitely many models}

\begin{prop}\label{prop:finmod}
Let $\cC$ be a pretopos with finite sorts such that $\Mod(\cC)$ has
finitely many isomorphism classes and finite sets of homomorphisms.
Let $\cX$ be a skeleton of $\Mod(\cC)$. Then the evaluation
$\mathrm{ev}\colon\cC\to\Fin^\cX$ is an equivalence.
\end{prop}
\begin{proof}
The evaluation is coherent and reflects the order of subobjects by
(R1); so it is faithful and every $\Sub_\cC(A)$ is finite. Choose a
finite set $S$ of objects such that two different homomorphisms
between models in $\cX$ differ on an object of $S$. For $M$ in $\cX$
let $\nu$ be the tuple of all elements of the sets $A(M)$, $A\in S$,
an element of $B(M)$ for a product $B$ of objects of $S$, and $S_M\le
B$ the meet of the finitely many subobjects containing $\nu$. For a
model $N$ and $\nu'\in S_M(N)$, every subobject of a product of
objects of $S$ which contains a tuple of entries of $\nu$ contains the
corresponding tuple of $\nu'$ (its pullback to $B$ is above $S_M$); by
(R2) there is a homomorphism $M\to N$ with $\nu\mapsto\nu'$, unique by
the choice of $S$. So $\mathrm{ev}(S_M)\cong\cX(M,-)$, $\nu$
corresponding to the identity. \emph{Full on subobjects:} for $W\le
\mathrm{ev}(A)$ and $w\in W(M)$ the Yoneda map $\cX(M,-)\to
\mathrm{ev}(A)$ of $w$ has as graph a subfunctor of $\mathrm{ev}(S_M
\times A)$; the subobject $\Gamma_w\le S_M\times A$, the meet of those
containing $(\nu,w)$, has $\mathrm{ev}(\Gamma_w)$ equal to this graph:
for $(\nu',w')\in\Gamma_w(N)$ apply (R2) to the family $(\nu,w)$, which
gives a homomorphism $\lambda$ with $\lambda\nu=\nu'$ and $\lambda
w=w'$. So the subfunctor generated by $w$ is $\mathrm{ev}$ of the
image of $\Gamma_w$ in $A$, and $W$ is a finite join of such. \emph{Full
and essentially surjective:} the graph of a morphism $\mathrm{ev}(A)
\to\mathrm{ev}(A')$ is $\mathrm{ev}(\Gamma)$, and $\Gamma$ is a graph
since $\mathrm{ev}$ reflects the order of subobjects; every functor
$\cX\to\Fin$ is a quotient of a finite coproduct of representables by
an equivalence relation, which comes from $\cC$.
\end{proof}

\begin{cor}[necessary condition]\label{cor:nec}
Let $\cC$ be a pretopos with finite sorts, $L$ a non-classical logic proving
linearity, and suppose that $\HH_L\cC$ is a pretopos.
Let $F$, $I$ be a filter and an ideal of $\Sub(1)$ such that the full
subcategory $\cX$ of $\Mod(\cC)$ of the models $M$ with $F\subseteq
\Sent(M)$, $\Sent(M)\cap I=\emptyset$ has finitely many isomorphism
classes and finite sets of homomorphisms. Then (a skeleton of) $\cX$
is a finite EI-category and every unit $W\to\pi^*\Lan_\pi W$, $W\colon
\cS_{n_L}(\cX)\to\Fin$, is injective.
\end{cor}
\begin{proof}
Theorem~\ref{thm:loc}, Proposition~\ref{prop:finmod} and (R5); $\cX$
is an EI-category by (R4).
\end{proof}


\section{Rigid models}

\begin{proof}[Proof of Theorem~\ref{thm:D}]
We may assume that $\cC$ is a pretopos. We use the notation of
\cite[\S5]{Finsorts}: for a finite language $T=(S,R)$, a finite set
$S$ of objects and a finite set $R$ of subobjects of products of
objects of $S$, $\bT_T$ is the theory of the coherent sequents in this
language valid in $\cC$, $\cX_T$ the finite category of its models and
$\cC_T\simeq\Fin^{\cX_T}$ its pretopos; $\cC$ is the directed colimit
of the $\cC_T$ \cite[Lemma~5.2]{Finsorts}. $\Mod(\cC)$ is an
EI-category. Let $S$ be a finite set of objects and $\cC_S$ the full
subcategory of the subobjects of the finite products of objects of
$S$. By Step~1 of the proof of \cite[Prop.~5.4]{Finsorts} every model
of $\cC_S$ is isomorphic to the restriction $M|_S$ of a model $M$ of
$\cC$, and every homomorphism $M|_S\to M'|_S$ is the restriction of a
homomorphism $M\to M'$; so every endomorphism of a model of $\cC_S$ is
an identity. Now we argue as in Step~3 there. Let $n_A$ be a bound for
the number of elements of $A\in S$ in the models, which exists by
compactness \cite[Lemma~4.1]{Finsorts}. For a finite set $R$
let $\mathcal E_R$ be the set of the pairs $(m,g)$ of a model $m$ of
$\bT_{(S,R)}$ whose sorts are sets $\{1,\dots,e_A\}$, $e_A\le n_A$, and
an endomorphism $g$ of $m$ which is not the identity on some sort.
These are finite sets, and restriction gives maps $\mathcal E_{R'}\to
\mathcal E_R$ for $R\subseteq R'$. If all were non-empty, an element
of the inverse limit would be a model of $\cC_S$ with an endomorphism
which is not an identity. So there is $R_0$ such that for $R\supseteq
R_0$ every endomorphism of a model of $\bT_{(S,R)}$ is an identity.
For these cofinally many $T$ the category $\cX_T$ is an EI-category in
which every automorphism is an identity, and $\HH_L\cC_T$ is a
pretopos by \cite[Cor.~6.5]{Strings}. Now (R3) applies.
\end{proof}

A category of models as in Theorem~\ref{thm:D} need not be a poset:
for $\cX=(0\to1\rightrightarrows2)$, with the two composites equal,
$\Fin^\cX$ is an example \cite[Ex.~6.6]{Strings}.

\section{Towers of finite categories}\label{sec:towers}

A \emph{tower} is an inverse system $(\cX_i)_{i\in I}$, $I$ directed,
of finite EI-categories with the same finite set of objects and
functors $q_{i'i}\colon\cX_{i'}\to\cX_i$, $i'\ge i$, which are the
identity on objects and full. Let $\cM$ be its limit: the same
objects, $\cM(x,y)=\varprojlim\cX_i(x,y)$, a profinite set; $q_i\colon
\cM\to\cX_i$ the projections; and $\cC=\varinjlim\Fin^{\cX_i}$ along
the functors $q^*$, a pretopos. Strings, $\cS_n(\cM)$, one-step
arrows, $V_t$, $K_t$, $G_s$, $H_s$ are defined for $\cM$ as for a
finite EI-category.

\begin{lem}\label{lem:tower}
\begin{enumerate}
\item[(a)] The functors $q_{i'i}$ and $q_i$ are surjective on arrows
and reflect isomorphisms, and $\cM$ is an EI-category.
\item[(b)] Every string of
$\cX_i$ is the image of a string of $\cM$; and for a string $\tilde t$
of $\cM$, every one-step arrow $\bar\beta\colon\bar t'\to q_i\tilde t$
of $\cX_i$ is $q_i\beta$ for a one-step arrow $\beta\colon\tilde t'\to
\tilde t$ of $\cM$ (with $q_i\tilde t'=\bar t'$). The same holds for
$q_{i'i}$.
\item[(c)] For strings $\tilde t,\tilde s$ of $\cM$:
$K_{\tilde t}=\bigcap_iq_i^{-1}(K_{q_i\tilde t})$ in $V_{\tilde t}$,
where $K_{q_i\tilde t}=\cS_n(\cX_i)(q_i\tilde t,q_i\tilde s)$, and for
every $i$ there is $i'\ge i$ with $q_{i'i}(K_{q_{i'}\tilde t})=
q_i(K_{\tilde t})$.
\end{enumerate}
\end{lem}
\begin{proof}
(a) An inverse limit of non-empty finite sets over a directed set is
non-empty, so $q_i$ is surjective on arrows. Let $q(f)$ be invertible,
$f\colon x\to y$; lift its inverse to $g\colon y\to x$; then $gf$ and
$fg$ are endomorphisms, invertible in an EI-category, so $f$ is
invertible. An endomorphism of $\cM$ has invertible images in all
$\cX_i$, and the inverses are compatible. (b) Non-invertible arrows
lift to non-invertible arrows, so strings lift. Let $\bar\beta=
\bar\iota_0\bar b$, with $\bar b$ the first arrow of $\bar t'$ and
$\bar\iota$ an isomorphism of strings from $\bar t'_{\ge1}$ onto
$q_i\tilde t$. Lift the $\bar\iota_k$ to isomorphisms $\iota_k$ of
$\cM$, let $\tilde t'_{\ge1}$ be the string with the arrows
$\iota_k^{-1}\tilde t_k\iota_{k-1}$, which lie over the arrows of
$\bar t'_{\ge1}$, and lift $\bar b$ to $b$; then $\tilde t'$ is a
string of $\cM$ over $\bar t'$, $\iota$ is an isomorphism of strings
$\tilde t'_{\ge1}\to\tilde t$, and $\beta=\iota_0b$ is a one-step arrow
over $\bar\beta$. (c) An arrow $g\in V_{\tilde t}$ lies in
$K_{\tilde t}$ iff $g=\iota_0p$ for an isomorphism of strings $\iota$
from a final segment of $\tilde t$ onto $\tilde s$; the index $j$ of
the segment is determined by the length of $\tilde s$; the pairs
$(g,\iota)$ at the stages form an inverse system of finite sets, so
$g$ with all $q_i(g)\in K_{q_i\tilde t}$ lifts to a compatible family
$\iota$. The second statement: the sets $q_{i'i}(K_{q_{i'}\tilde t})$
decrease in the finite set $V_{q_i\tilde t}$, and an element of all of
them is the image of an element of $\bigcap_{i'}q_{i'}^{-1}(K_{q_{i'}
\tilde t})$ by compactness.
\end{proof}


\begin{lem}\label{lem:trans}
Let $q\colon\cX'\to\cX$ be a functor between finite EI-categories
which reflects isomorphisms, and $L$ a logic proving linearity. Then
$q$ induces a functor $q\colon\cS_n(\cX')\to\cS_n(\cX)$, and under the
equivalences of \emph{(R5)} the functor $\PP\HH_L(q^*)$ is $W\mapsto
Wq$.
\end{lem}
\begin{proof}
$q$ sends non-invertible arrows to non-invertible arrows, hence
strings to strings, and isomorphisms of strings to isomorphisms of
strings. For a string $s'$ of $\cX'$, the composite of $q^*$ with the
Kripke model $s'$ of $\Fin^{\cX'}$ is the Kripke model $qs'$ of
$\Fin^\cX$, as the evaluation at an object $x'$ composed with $q^*$ is
the evaluation at $qx'$; a Heyting functor from $\HH_L$ to the functors
on a chain is determined up to isomorphism by its restriction, so
$\hat G_{s'}\HH_L(q^*)\cong\hat G_{qs'}$. Hence the two functors agree
on $\HH_L\Fin^\cX$, compatibly with the arrows of $\cS_n$, and both
are coherent functors on its pretopos completion.
\end{proof}

\begin{lem}[the eventual criterion]\label{lem:eventual}
Let $(\cX_i)$ be a tower, $\cC=\varinjlim\Fin^{\cX_i}$ and $L$ a logic
proving linearity. Then $\HH_L\cC$ is a pretopos if and only if for
every $i$ and every $W\colon\cS_n(\cX_i)\to\Fin$ there is $i'\ge i$
such that the unit of $Wq_{i'i}$ is injective.
\end{lem}
\begin{proof}
By (R3) $\HH_L\cC$ is the colimit of the $\HH_L\Fin^{\cX_i}$, which by
(R5) and Lemma~\ref{lem:trans} are the full subcategories of the
$\Fin^{\cS_n(\cX_i)}$ of the functors with injective unit, with the
transitions $W\mapsto Wq$.

Let the condition hold. An equivalence relation in $\HH_L\cC$ is
represented by one at a stage $i$, on an object $Z$; let $W=Z/E$ in
$\Fin^{\cS_n(\cX_i)}$. For $i'$ as in the condition, $Wq_{i'i}$ is an
object of $\HH_L\Fin^{\cX_{i'}}$ and the quotient of $Zq$ by $Eq$
there and at all later stages, the transitions being restrictions of
coherent functors between pretoposes; so it is an effective quotient
in the colimit, stable under pullback. The same argument gives the
finite coproducts, which are disjoint. So $\HH_L\cC$ is a pretopos.

Conversely let $\HH_L\cC$ be a pretopos and $W$ a functor at the stage
$i$. It is the quotient of a finite coproduct of objects $Z_k$ of
$\HH_L\Fin^{\cX_i}$ (for instance representable functors, which are
subfunctors of functors $\pi^*A$) by an equivalence relation given by
subobjects $E_{kl}\le Z_k\times Z_l$. In the pretopos $\HH_L\cC$ let
$Q$ be the quotient of the coproduct of the $Z_k$ by this relation.
The object $Q$ and the morphisms $Z_k\to Q$ are represented at a stage
$i'\ge i$, where they are jointly covering with the kernel relations
$E_{kl}q$, these being finitely many conditions on subobjects. In the
pretopos $\Fin^{\cS_n(\cX_{i'})}$ this identifies the representative
of $Q$, an object of $\HH_L\Fin^{\cX_{i'}}$, with $Wq_{i'i}$; so the
unit of $Wq_{i'i}$ is injective.
\end{proof}


\section{Profinite leg trees}\label{sec:trees}

Let $(\cX_i)$ be a tower with limit $\cM$.

\begin{lem}\label{lem:mono}
\begin{enumerate}
\item[(a)] If there is a morphism of leg trees $T\to T''$, then
$E_T\subseteq E_{T''}$ and $F_T\subseteq F_{T''}$.
\item[(b)] Let $q\colon\cX'\to\cX$ be a functor between finite
EI-categories which reflects isomorphisms, and $T'$ a leg tree of
$\cX'$ at $t'$ over $s'$. Applying $q$ to the labels gives a leg tree
$qT'$ of $\cX$ at $qt'$, and $q(E_{T'})\subseteq E_{qT'}$,
$q(F_{T'})\subseteq F_{qT'}$.
\item[(c)] For $W$, $s$ and $a\in W(s)$ let $\mathrm{Max}(s,a)$ be the
leg tree at $s$ defined by recursion: the children of the root are the
pairs $(\beta,y)$ of a one-step arrow $\beta\colon t'\to s$ and $y\in
W(t')$ with $W(\beta)y=a$, with label $\beta$ and subtree
$\mathrm{Max}(t',y)$. It is finite and realised in $W$ with $a$ at the
root, and every leg tree at $s$ realised in $W$ with $a$ at the root
has a morphism to it. Hence the unit of $W$ is injective at $s$ iff
for all $a\in W(s)$ and $\Theta\in H_s$ with $\id\mathrel{E_{
\mathrm{Max}(s,a)}}\Theta$ we have $W(\Theta)a=a$.
\end{enumerate}
\end{lem}
\begin{proof}
(a) By induction: a child $(\beta,T_1)$ of the root of $T$ is sent to
a child $(\beta,T_1'')$ of the root of $T''$ with a morphism $T_1\to
T_1''$; so $\beta^*E_{T_1}\subseteq\beta^*E_{T_1''}\subseteq E_{T''}$.
(b) $q$ sends strings to strings, arrows of $\cS_n$ to arrows of
$\cS_n$ and one-step arrows to one-step arrows, and maps $V_{t'}$ to
$V_{qt'}$ and $K_{t'}$ to $K_{qt'}$; the inclusions follow by induction
since $q(g\beta)=q(g)q(\beta)$. (c) $\cS_n(\cX)$ is a finite category
and the strings get longer, so the tree is finite; the elements $y$
are a realisation; a realised tree maps to it by sending a node with
element $y_\nu$ and label $\beta$ to the child $(\beta,y_\nu)$,
recursively. The last statement follows from (R6) and (a): if the unit
identifies $a$ and $b=W(\Theta)a$ through $T$, then $\id\mathrel{
E_{\mathrm{Max}(s,a)}}\Theta$; conversely $\mathrm{Max}(s,a)$ is
itself a realised tree.
\end{proof}

Let $\tilde t$ be a string of $\cM$ over $\tilde s$, $i\in I$ and $T$
a leg tree of $\cX_i$ at $q_i\tilde t$. For $i'\ge i$ the \emph{pullback}
$T^{i'}$ of $T$ to $\tilde t$ is the leg tree of $\cX_{i'}$ at $q_{i'}
\tilde t$ defined by recursion: the children of its root are the pairs
$(\beta',\nu)$ of a child $\nu=(\beta,T_1)$ of the root of $T$ and a
one-step arrow $\beta'\colon t'\to q_{i'}\tilde t$ of $\cX_{i'}$ with
$q\beta'=\beta$ (so $qt'$ is the string of $\nu$), with label $\beta'$
and subtree the pullback of $T_1$ to (a string of $\cM$ over) $t'$;
the relations $E$, $F$ of a pullback do not depend on the choice of
the strings of $\cM$, being defined in $\cX_{i'}$. By
Lemma~\ref{lem:mono}(b) the relations
\[
\mathcal E_{i'}=q_{i'}^{-1}(E_{T^{i'}})\ \text{on }V_{\tilde t},\qquad
\mathcal F_{i'}=\{(\kappa,\kappa')\in K_{\tilde t}^2\mid q_{i'}\kappa
\mathrel{F_{T^{i'}}}q_{i'}\kappa'\}
\]
decrease with $i'$ (the image of $T^{i''}$ in $\cX_{i'}$ has a
morphism to $T^{i'}$, and Lemma~\ref{lem:mono}(a)). The pair
$\mathcal T=(i,T)$ is called a \emph{profinite leg tree} at $\tilde
t$, and
\[
\bar E_{\mathcal T}=\bigcap_{i'\ge i}\mathcal E_{i'},\qquad
\bar F_{\mathcal T}=\bigcap_{i'\ge i}\mathcal F_{i'}\subseteq\bar
E_{\mathcal T}\cap(K_{\tilde t}\times K_{\tilde t}).
\]
$\mathcal T$ is an \emph{obstruction} if $\bar F_{\mathcal T}\ne\bar
E_{\mathcal T}\cap(K_{\tilde t}\times K_{\tilde t})$.

\begin{thm}\label{thm:proftree}
Let $(\cX_i)$ be a tower with limit $\cM$, $\cC=\varinjlim\Fin^{\cX_i}$
and $L$ a logic proving linearity, $n=n_L$. Then $\HH_L\cC$
is a pretopos if and only if for every string $\tilde s$ of $\cM$ no
profinite leg tree at $\tilde s$ (over $\tilde s$) is an obstruction.
\end{thm}
\begin{proof}
By Lemma~\ref{lem:eventual}, $\HH_L\cC$ is a
pretopos iff for every $i$ and $W\colon\cS_n(\cX_i)\to\Fin$ there is
$i'\ge i$ such that the unit of $Wq_{i'i}$ is injective. At the root
of a leg tree at $\tilde s$ the relations are invariant under left
multiplication, at every stage and so in the limit; so it suffices to
consider the classes of the identity.

Fix $i$ and $W$. For a string $\tilde s$ of $\cM$ and $a\in W(q_i
\tilde s)$ let $T=\mathrm{Max}(q_i\tilde s,a)$ for $W$. For $i'\ge i$
the tree $\mathrm{Max}(q_{i'}\tilde s,a)$ for $Wq_{i'i}$ is the
pullback $T^{i'}$: its children at the root are the $(\beta',y)$ with
$W(q\beta')y=a$, that is with $(q\beta',y)$ a child of the root of
$T$, and recursively. Let
\[
B_{i'}(\tilde s,a)=q_{i'i}\bigl(\{\Theta'\in H_{q_{i'}\tilde s}\mid
\id\mathrel{E_{T^{i'}}}\Theta'\}\bigr)\subseteq G_{q_i\tilde s}.
\]
By Lemma~\ref{lem:mono}(c), the unit of $Wq_{i'i}$ is injective iff
$W(\theta)a=a$ for all $\tilde s$, $a$ and $\theta\in B_{i'}(\tilde
s,a)$ (every string of $\cX_{i'}$ is some $q_{i'}\tilde s$). The sets
$B_{i'}(\tilde s,a)$ decrease with $i'$ and depend only on $q_{i'}
\tilde s$ and $a$; by Lemma~\ref{lem:tower}(c) (for $\tilde t=\tilde
s$) and compactness their intersection is
\[
B_\infty(\tilde s,a)=q_i\bigl(\{\Theta\in H_{\tilde s}\mid\id\mathrel{
\bar E_{\mathcal T}}\Theta\}\bigr),\qquad\mathcal T=(i,T).
\]
The strings of $\cM$ of a given shape form a profinite space, a closed
subspace of a product of sets of arrows, on which $q_{i'}$ is locally
constant. If $B_\infty(\tilde s,a)\subseteq\mathrm{Stab}(a)$ for all
$\tilde s$, $a$, then for each $\tilde s$ there is $i'$ with $B_{i'}
(\tilde s,a)\subseteq\mathrm{Stab}(a)$ for all $a$, and this holds for
all strings with the same image in $\cX_{i'}$, an open set, and for
all larger $i'$; finitely many of these open sets cover, so one $i'$
serves for all $\tilde s$. Hence:
\begin{equation}\label{eq:W}
\text{$W$ is repaired at a later stage iff $B_\infty(\tilde
s,a)\subseteq\mathrm{Stab}_W(a)$ for all $\tilde s$ and $a$.}
\end{equation}

\emph{Sufficiency.} Suppose no profinite leg tree at a string $\tilde
s$ is an obstruction, and let $\Theta\in H_{\tilde s}$ with $\id
\mathrel{\bar E_{\mathcal T}}\Theta$, $\mathcal T=(i,T)$ as above.
Then $\id\mathrel{\bar F_{\mathcal T}}\Theta$, in particular $\id
\mathrel{F_T}q_i\Theta$ in $\cX_i$ ($T^i=T$ up to a morphism in both
directions), and $W(q_i\Theta)a=a$ by (R6), $T$ being realised in $W$.
So \eqref{eq:W} holds for every $W$.

\emph{Necessity.} Let $\mathcal T=(i,T)$ be a profinite leg tree at
$\tilde s$ and $\Theta\in H_{\tilde s}$ with $\id\mathrel{\bar
E_{\mathcal T}}\Theta$; we show $\id\mathrel{\mathcal F_{i'}}\Theta$
for every $i'\ge i$. Since $(i',T^{i'})$ has the same pullbacks as
$\mathcal T$ at the stages above $i'$, up to morphisms in both
directions, we may assume $i'=i$. Let $W=W_T$ over $\cX_i$ as in (R7),
with $a$ at the root. $T$ is realised in $W$, so it has a morphism to
$\mathrm{Max}(q_i\tilde s,a)$, and the same holds for the pullbacks;
by Lemma~\ref{lem:mono}(a) $q_i\Theta\in B_\infty(\tilde s,a)$. As $W$
is repaired at a later stage, \eqref{eq:W} gives $W(q_i\Theta)a=a$,
that is $\id\mathrel{F_T}q_i\Theta$.
\end{proof}

For a set $\rho_\lambda$, $\lambda\in\Lambda$, of relations on a
profinite set $V$ let $\widehat\bigvee_\lambda\rho_\lambda$ be the
intersection of all open equivalence relations on $V$ (those with
finitely many classes, all open) containing all $\rho_\lambda$.

\begin{prop}[recursion]\label{prop:rec}
Let $\mathcal T=(i,T)$ be a profinite leg tree at $\tilde t$. If the
root of $T$ has no children, $\bar E_{\mathcal T}$ and $\bar
F_{\mathcal T}$ are the equality. Otherwise
\[
\bar E_{\mathcal T}=\widehat\bigvee\beta^*\bar E_{\mathcal T_{\beta,\nu}},
\qquad\bar F_{\mathcal T}=\widehat\bigvee\beta^*\bar F_{\mathcal T_{\beta,\nu}},
\]
over the pairs $(\beta,\nu)$ of a child $\nu=(q_i\beta,T_1)$ of the
root of $T$ and a one-step arrow $\beta\colon\tilde t'\to\tilde t$ of
$\cM$ over its label, where $\mathcal T_{\beta,\nu}=(i,T_1)$ at $\tilde t'$, the
joins being formed on $V_{\tilde t}$ and on $K_{\tilde t}$.
\end{prop}
\begin{proof}
If the root has no children, $\mathcal E_{i'}$ is the relation
$q_{i'}g=q_{i'}g'$, with intersection the equality.

\emph{$\supseteq$.} $\mathcal E_{i'}$ is an open equivalence relation,
and it contains $\beta^*q_{i'}^{-1}(E_{T_1^{i'}})\supseteq\beta^*\bar
E_{\mathcal T_{\beta,\nu}}$ since $(q_{i'}\beta,\nu)$ is a child of the root
of $T^{i'}$. The same for $F$.

\emph{$\subseteq$ for $E$.} Let $\mathcal E$ be an open equivalence
relation on $V_{\tilde t}$ containing all $\beta^*\bar E_{\mathcal T_{\beta,\nu}}$. For each $\beta$ the closed relations $\beta^*q_{i'}^{-1}
(E_{T_1^{i'}})$ decrease to $\beta^*\bar E_{\mathcal T_{\beta,\nu}}\subseteq
\mathcal E$, so one of them is contained in $\mathcal E$; it depends
only on $q_{i'}\beta$ and the child. The pairs $(\beta,\nu)$ form a
profinite space (the inverse limit of the finite sets of the children
of the roots of the $T^{i'}$, by Lemma~\ref{lem:tower}(b)), so there is
$i'$, which we take so large that $\mathcal E=q_{i'}^{-1}(\mathcal
E')$ for an equivalence relation $\mathcal E'$ on $V_{q_{i'}\tilde
t}$, with $\beta^*q_{i'}^{-1}(E_{T_1^{i'}})\subseteq\mathcal E$ for
all $(\beta,\nu)$. Every child of the root of $T^{i'}$ is $(q_{i'}
\beta,\nu)$ for some $\beta$, and $q_{i'}$ is surjective on $V$; so
$(q_{i'}\beta)^*E_{T_1^{i'}}\subseteq\mathcal E'$ for all children,
$E_{T^{i'}}\subseteq\mathcal E'$ and $\bar E_{\mathcal T}\subseteq
\mathcal E_{i'}\subseteq\mathcal E$.

\emph{$\subseteq$ for $F$.} Let $\mathcal E$ be an open equivalence
relation on $K_{\tilde t}$ containing all $\beta^*\bar F_{\mathcal T_{\beta,\nu}}$. As before there is $i'$ such that $\mathcal E$ is the
inverse image of a relation $\mathcal E'$ on $q_{i'}(K_{\tilde t})$
and $(\kappa\beta,\kappa'\beta)\in\mathcal F^{\mathcal T_{\beta,\nu}}_{i'}$
implies $(\kappa,\kappa')\in\mathcal E$, for $\kappa,\kappa'\in
K_{\tilde t}$ and all $(\beta,\nu)$. By Lemma~\ref{lem:tower}(c) choose
$i''\ge i'$ with $q_{i''i'}(K_{q_{i''}\tilde t})=q_{i'}(K_{\tilde
t})$. Let $(\kappa,\kappa')\in\mathcal F_{i''}$: there is a chain
$q_{i''}\kappa=\bar\kappa_0,\dots,\bar\kappa_r=q_{i''}\kappa'$ in
$K_{q_{i''}\tilde t}$ with $\bar\kappa_l\bar\beta_l\mathrel{F}
\bar\kappa_{l+1}\bar\beta_l$ for children $(\bar\beta_l,\nu_l)$ of the
root of $T^{i''}$ and $F$ the relation of the subtree. Its image in
$\cX_{i'}$ is a chain in $q_{i'}(K_{\tilde t})$ whose consecutive
members are related in the same way by children of the root of
$T^{i'}$ (Lemma~\ref{lem:mono}(a),(b)). Lift its members to $\kappa_l
\in K_{\tilde t}$ and the children to $(\beta_l,\nu_l)$; then
$(\kappa_l,\kappa_{l+1})\in\mathcal E$, and $(\kappa,\kappa_0)$,
$(\kappa_r,\kappa')\in\mathcal E$ as these have the same images in
$\cX_{i'}$. So $\bar F_{\mathcal T}\subseteq\mathcal F_{i''}\subseteq
\mathcal E$.
\end{proof}

\begin{rem}\label{rem:groups}
At the root $\tilde s$ the relations are invariant under left
multiplication by $G_{\tilde s}$ resp.\ $H_{\tilde s}$, so $\bar
E_{\mathcal T}$ is $g^{-1}g'\in S_{\mathcal T}$ and $\bar F_{\mathcal
T}$ is $h^{-1}h'\in R_{\mathcal T}$ for closed subgroups $R_{\mathcal
T}\le H_{\tilde s}\cap S_{\mathcal T}$ of the profinite group
$G_{\tilde s}$, and the condition of Theorem~\ref{thm:proftree} is
$R_{\mathcal T}=H_{\tilde s}\cap S_{\mathcal T}$.
\end{rem}

\begin{prop}[the towers are general]\label{prop:general}
Let $\cC$ be a pretopos with finite sorts such that $\Mod(\cC)$ is an
EI-category and $\Sub(1)$ is finite. Then $\cC\simeq\varinjlim
\Fin^{\cX_S}$ for a tower $(\cX_S)$ indexed by the finite sets $S$ of
objects of $\cC$, with limit equivalent to $\Mod(\cC)$.
\end{prop}
\begin{proof}
For a finite set $S$ of objects let $\cC_S$ be the full subcategory of
the subobjects of the finite products of objects of $S$, a coherent
category with finite sorts and with the same sentences as $\cC$. By
Step~1 of the proof of \cite[Prop.~5.4]{Finsorts}, which uses that
$\Mod(\cC)$ is an EI-category, every model of $\cC_S$ is isomorphic to the restriction
$M|_S$ of a model of $\cC$, and by (R2) every homomorphism $M|_S\to
M'|_S$ is the restriction of a homomorphism $M\to M'$. Choose one
model $M_P$ for each of the finitely many prime filters $P$ of
$\Sub(1)$ (Proposition~\ref{prop:spec}), and let $\cX_S$ have these as
objects and $\cX_S(P,Q)$ the set of the restrictions to $\cC_S$ of the
homomorphisms $M_P\to M_Q$, a finite set. It is a finite EI-category;
the models $M_P|_S$ are pairwise non-isomorphic, having different
sentences; so $\cX_S$ is a skeleton of $\Mod(\cC_S)=\Mod(\PP\cC_S)$,
and $\PP\cC_S\simeq\Fin^{\cX_S}$ by Proposition~\ref{prop:finmod}. For
$S\subseteq S'$ restriction is a functor $\cX_{S'}\to\cX_S$, the
identity on objects and full. The limit has the hom-sets $\varprojlim_S
\cX_S(P,Q)$, the homomorphisms $M_P\to M_Q$. Finally $\cC$ is the
colimit of the $\PP\cC_S$: the functors $\PP\cC_S\to\cC$ reflect the
order of subobjects, every object $A$ of $\cC$ lies in $\cC_{\{A\}}$,
and a morphism between objects of $\PP\cC_S$, quotients of coproducts
of objects of $\cC_S$, is given by subobjects of products of objects
of $S$, which are in $\cC_S$; so the functors $\PP\cC_S\to\cC$ are
full and faithful and jointly essentially surjective.
\end{proof}


\begin{cor}\label{cor:A}
Let $\cC$ be a pretopos with finite sorts such that $\Sub(1)$ is
finite, and $L$ a non-classical logic proving linearity. Then $c^L_\cC$ is an equivalence if and only if $\Mod(\cC)$ is
an EI-category and no profinite leg tree of the tower of
Proposition~\ref{prop:general}, at a string of models over itself, is
an obstruction.
\end{cor}
\begin{proof}
If $c^L_\cC$ is an equivalence, $\Mod(\cC)$ is an EI-category by (R4).
If $\Mod(\cC)$ is an EI-category, Proposition~\ref{prop:general} and
Theorem~\ref{thm:proftree} apply.
\end{proof}

This is Theorem~\ref{thm:A}. By Proposition~\ref{prop:rec} and
Remark~\ref{rem:groups} the subgroups $R_{\mathcal T}\le H_s\cap
S_{\mathcal T}$ are computed in $\Mod(\cC)$ from the profinite sets of
homomorphisms, once the legs of $\mathcal T$ at each node are given.

\begin{ex}\label{ex:product}
Let $X$ be a non-empty finite poset and $\cC$ the pretopos of the
theory of \cite[Ex.~7.3]{Finsorts}: propositions $u_x$, $x\in X$, with
axioms saying that the true ones form the down-set of an element of
$X$, and countably many sorts with exactly two elements, each with a
relation $\ne$ which is the complement of the equality.
Its category of models is equivalent to $X\times\Gamma$ for the
profinite group $\Gamma=(\mathbb Z/2)^{\mathbb N}$, regarded as a
category with one object, and $\Sub(1)$ is the lattice of the up-sets
of $X$. The tower of the categories $X\times(\mathbb Z/2)^e$ has this
limit, and the same colimit of pretoposes. In $X\times(\mathbb
Z/2)^e$ the arrows $t\to s$ of $\cS_n$ are all the arrows $t_0\to
s_0$ if the chain of $s$ is a final segment of that of $t$, and there
are none otherwise \cite[proof of Cor.~6.2]{Strings}; so $K_t=V_t$ for
every string $t$ over $s$, the two recursions defining $E_T$ and
$F_T$ coincide at every stage, and no profinite leg tree is an
obstruction. By Theorem~\ref{thm:proftree} the comparison is an
equivalence for every $L$ proving linearity, in accordance with
\cite[Ex.~7.3]{Finsorts}.
\end{ex}

\section{A family with profinite abelian groups}\label{sec:three}

Let $G$ be a profinite abelian group and $A,T_1,\dots,T_m$ closed
subgroups. For an open subgroup $N$ let $\cX_N$ be the category with
the objects $0_1,\dots,0_m,1,2$, with $\Aut(2)=(A+N)/N$ and trivial
automorphism groups otherwise, and
\[
\cX_N(1,2)=G/N,\qquad\cX_N(0_j,1)=\cX_N(0_j,2)=G/(T_j+N),
\]
composition being the addition; this is the finite EI-category of
\cite[Prop.~6.11]{Strings} for the group $G/N$ and the images of the
subgroups. For $N'\subseteq N$ the quotient maps are a functor
$\cX_{N'}\to\cX_N$, the identity on objects and full. So the $\cX_N$
are a tower; let $\cC(G;A;T_\bullet)=\varinjlim\Fin^{\cX_N}$. The limit
$\cM$ has $\Aut(2)=A$, $\cM(1,2)=G$ and $\cM(0_j,1)=\cM(0_j,2)=G/T_j$,
the subgroups being closed. The pretopos $\cC(G;A;T_\bullet)$ has
finite sorts, its models are the objects of $\cM$ and the
homomorphisms are the arrows: a model restricts to a model of each
$\Fin^{\cX_N}$, an object of $\cX_N$, the same for all $N$, and a
homomorphism to a compatible family of arrows. Its sentences are the
up-sets of the poset $0_j<1<2$.

\begin{proof}[Proof of Theorem~\ref{thm:B}, first part]
The strings have at most three objects. For a string $s$ with first
object $0_j$ or $1$ the group $G_s$ is trivial, so no profinite leg
tree at $s$ is an obstruction (Remark~\ref{rem:groups}). Let $s=(2)$,
with $G_s=H_s=A$. The one-step arrows into $s$ are the arrows
$\kappa\colon t_\gamma\to s$ from the strings $t_\gamma=(1\xrightarrow
{\gamma}2)$, $\gamma\in G$, with $\kappa\in A+\gamma$, and the arrows
$c\colon(0_j\xrightarrow{d}2)\to s$ with $c\in A+d\subseteq G/T_j$. The one-step arrows into
$t_\gamma$ are the arrows $x\colon(0_j\xrightarrow{x}1\xrightarrow
{\gamma'}2)\to t_\gamma$ with $\gamma'\in A+\gamma$, as $\Aut(1)$ is
trivial; these strings have no one-step arrows into them. For
$t=t_\gamma$ over $s$ we have $V_t=G$ and $K_t=A+\gamma$, and for
$u=(0_j\to1\to2)$ the set $V_u=G/T_j$.

Let $\mathcal T=(N,T)$ be a profinite leg tree at $s$. We compute the
relations by Proposition~\ref{prop:rec}. At a node with a string $u$
they are the equality. At a node $\nu$ with the string $t_\gamma$,
with $J_\nu$ the set of the $j$ such that $\nu$ has a child with first
object $0_j$: the leg $x$ gives on $V_t=G$ the relation $g+x=g'+x$ in
$G/T_j$, that is $g-g'\in T_j$, and on $K_t$ the relation
$\kappa-\kappa'\in A\cap T_j$. An open equivalence relation on $G$
containing the cosets of the $T_j$, $j\in J_\nu$, contains the cosets
of the closed subgroup $T_{J_\nu}=\sum_{j\in J_\nu}T_j$, which is the
intersection of the open subgroups containing it; so $\bar E$ at $\nu$
is $g-g'\in T_{J_\nu}$, and likewise $\bar F$ at $\nu$ is
$\kappa-\kappa'\in\sum_{j\in J_\nu}(A\cap T_j)$. At the root, the leg
$\kappa$ into a child $\nu$ of this kind gives on $A$ the relations
$a-a'\in A\cap T_{J_\nu}$ for $\bar E$ and $a-a'\in\sum_{j\in J_\nu}
(A\cap T_j)$ for $\bar F$, and a leg $c$ from a string $(0_j\to2)$ gives
$a+c=a'+c$ in $G/T_j$, that is $a-a'\in A\cap T_j$, for both. Hence
\[
S_{\mathcal T}=\sum_\nu(A\cap T_{J_\nu})+\sum_c(A\cap T_{j_c}),\qquad
R_{\mathcal T}=\sum_\nu\sum_{j\in J_\nu}(A\cap T_j)+\sum_c(A\cap
T_{j_c}).
\]
If \eqref{eq:dist} holds, these are equal. Conversely, for
$J\subseteq\{1,\dots,m\}$ let $T$ be the leg tree of $\cX_N$ at $s$
with one child $(0\colon t_0\to s)$ of the root, which has the
children $(0\colon(0_j\xrightarrow{0}1\xrightarrow{0}2)\to t_0)$,
$j\in J$; then $S_{\mathcal T}=A\cap T_J$ and $R_{\mathcal T}=\sum_{j
\in J}(A\cap T_j)$. By Theorem~\ref{thm:proftree}, $\HH_L\cC$ is a
pretopos iff \eqref{eq:dist} holds.
\end{proof}

\begin{prop}\label{prop:cofinalN}
Let $G$, $A$, $T_1,\dots,T_m$ be as above and
suppose \eqref{eq:dist}.
Then every open subgroup of $G$ contains an open subgroup $N$ with
\[
(A+N)\cap\sum_{j\in J}(T_j+N)=\sum_{j\in J}\bigl((A+N)\cap(T_j+N)\bigr)
\qquad\text{for all }J,
\]
that is, such that $\HH_L\Fin^{\cX_N}$ is a pretopos for $L$ proving
linearity with $n_L\ge3$
\cite[Prop.~6.11]{Strings}.
\end{prop}
\begin{proof}
Let $\Gamma$ be the Pontryagin dual of $G$, a discrete torsion abelian
group. $B\mapsto B^\perp$ is an order-reversing bijection from the
closed subgroups of $G$ to the subgroups of $\Gamma$ which exchanges
sums and intersections, and the open subgroups correspond to the
finite ones. Let $a=A^\perp$, $t_j=T_j^\perp$, $t_J=\bigcap_{j\in J}
t_j$, and $\pi\colon\Gamma\to\Gamma/a$. The hypothesis says $a+t_J=
\bigcap_{j\in J}(a+t_j)$, that is $\pi(t_J)=\bigcap_{j\in J}\pi(t_j)$,
for all non-empty $J$ (for empty $J$ both conditions are trivial). For
$N$ with $F=N^\perp$ the required condition says $(a\cap F)+(t_J\cap
F)=\bigcap_{j\in J}((a\cap F)+(t_j\cap F))$; as $(a\cap F)+(t\cap F)=
F\cap\pi^{-1}\pi(t\cap F)$ for a subgroup $t$, this is
\begin{equation}\label{eq:stage}
\pi(t_J\cap F)=\bigcap_{j\in J}\pi(t_j\cap F)\qquad\text{for all
non-empty }J.
\end{equation}
Let $F_0$ be a finite subgroup of $\Gamma$. For each non-empty $J$ the
group $\pi(t_J)\cap\pi(F_0)$ is finite; choose a finitely generated,
hence finite, subgroup $\tau_J\le t_J$ with $\pi(\tau_J)=\pi(t_J)\cap
\pi(F_0)$, and let $F=F_0+\sum_J\tau_J$. Then $\pi(F)=\pi(F_0)$ and
\[
\begin{aligned}
\bigcap_{j\in J}\pi(t_j\cap F)&\subseteq\bigcap_{j\in J}\pi(t_j)\cap
\pi(F)=\pi(t_J)\cap\pi(F_0)=\pi(\tau_J)\\
&\subseteq\pi(t_J\cap F)
\subseteq\bigcap_{j\in J}\pi(t_j\cap F),
\end{aligned}
\]
so \eqref{eq:stage} holds for the finite group $F\supseteq F_0$.
\end{proof}


\begin{proof}[Proof of Theorem~\ref{thm:B}, second part]
This is Proposition~\ref{prop:cofinalN}.
\end{proof}

\begin{ex}\label{ex:Zp}
Let $G=\mathbb Z_p^3$, $T_1=\mathbb Z_pe_1$, $T_2=\mathbb Z_pe_2$ and
$A=\mathbb Z_p\alpha$, $\alpha=(1,1,p)$. Then $A\cap(T_1+T_2)=0$ by
the third coordinate, so \eqref{eq:dist} holds and $c^L_\cC$ is an
equivalence for every $L$ proving linearity (for
$n_L\le2$ by \cite[Thm.~1.1]{Finsorts}).

(a) For $N=p^kG$, $k\ge1$, the quotient does not satisfy
\eqref{eq:dist}: $(A+N)\cap(T_1+T_2+N)$ contains $p^{k-1}\alpha\notin
N$, as $p^{k-1}\alpha\equiv p^{k-1}(e_1+e_2)$ modulo $N$, while
$(A+N)\cap(T_j+N)=N$: $a\alpha\equiv be_1$ modulo $p^k$ forces
$a\equiv0$ by the second coordinate, and likewise for $e_2$ by the
first. So for $n_L\ge3$ the comparison is not an equivalence for any
of the pretoposes $\Fin^{\cX_{p^kG}}$ \cite[Prop.~6.11]{Strings},
whose union is $\cC$.

(b) In the basis $\alpha,e_2,e_3$ of $G$, where $e_1=\alpha-e_2-pe_3$,
let $N_{m,k}=p^m\mathbb Z_p\alpha+p^k\mathbb Z_pe_2+p^k\mathbb
Z_pe_3$ for $0\le m\le k-1$. Then $(A+N)\cap(T_2+N)=N$,
$(A+N)\cap(T_1+N)=p^k\mathbb Z_p\alpha+N$ and
$(A+N)\cap(T_1+T_2+N)=p^{k-1}\mathbb Z_p\alpha+N$, and all three are
equal to $N$ as $m\le k-1$. As $N_{r,r+1}\subseteq p^rG$, the $\cX_{N_{m,k}}$ are a cofinal
family for which the comparison is an equivalence; and $p^kG\subseteq
N_{m,k}$, so the quotients of (a) refine those of (b).
\end{ex}

So the absence of obstructions is not monotone along the finite
approximations of a theory.

\section{Questions}

\begin{que}
Is there a criterion as in Theorem~\ref{thm:A} when $\Sub(1)$ is
infinite? The categories $\cX_S$ of Proposition~\ref{prop:general} are then
infinite, and for the finite languages of \cite[\S5]{Finsorts}, which
contain finitely many sentences, non-isomorphic models can have
isomorphic restrictions. By Theorem~\ref{thm:C} it is necessary that
the parts of $\Mod(\cC)$ cut out by a filter and an ideal with
finitely many prime filters between them satisfy the condition of
Theorem~\ref{thm:A}.
\end{que}

\begin{que}
In the situation of Theorem~\ref{thm:A}: if $c^L_\cC$ is an
equivalence, are there cofinally many finite approximations for which
it is? Theorem~\ref{thm:B} says yes for the family of
Section~\ref{sec:three} and the approximations $\cX_N$.
\end{que}

\section*{Declarations}
The author has no competing interests to declare.

\begin{thebibliography}{9}
\bibitem{Interm} R.~Shibusawa, \emph{Intermediate logics and
imaginaries: the translation into an intermediate logic commutes with
the pretopos completion of all finite propositional theories exactly
for the extensions of G\"odel--Dummett logic}, preprint (2026),
archived at DOI 10.5281/zenodo.24150610.
\bibitem{Strings} R.~Shibusawa, \emph{G\"odel--Dummett logic over a
finite category of models: the models of the translated theory are
strings of arrows, and a subgroup obstruction to commuting with the
imaginaries}, preprint (2026), archived at DOI
10.5281/zenodo.24327653.
\bibitem{Finsorts} R.~Shibusawa, \emph{Intermediate logics and the
imaginaries of a theory whose models are finite on every sort},
preprint (2026), archived at DOI 10.5281/zenodo.24356860.
\bibitem{Ye24} L.~Ye, \emph{Stack representation of finitely presented
Heyting pretoposes I}, preprint, arXiv:2402.13099.
\bibitem{Ye26} L.~Ye, \emph{Duality theory for categorical theories},
preprint, arXiv:2605.22071.
\end{thebibliography}

\end{document}
